% ── Section IV: IMPLEMENTATION ────────────────────────────────────────────────
% INSERT THIS ENTIRE BLOCK between the end of Section III (\subsection{Safety Validator})
% and before \section{Experimental Evaluation}
% ─────────────────────────────────────────────────────────────────────────────

\section{Implementation}

\subsection{Data Ingestion Pipeline}

The ingestion pipeline processes four modalities independently before
consolidating all documents into a single ChromaDB vector store.

\textit{Social media text.} Up to 500 tweets labelled ``Related and
informative'' were retained from the CrisisLex-T26 2013 Colorado Floods
subset~\cite{b21} after applying the dataset relevance filter. Tweet
timestamps were reconstructed from Twitter Snowflake identifiers by
right-shifting the 64-bit identifier by 22 bits and adding the Twitter
epoch offset of 1{,}288{,}834{,}974{,}657\,ms, producing UTC timestamps not
available in the raw dataset metadata.

\textit{Satellite imagery descriptions.} Fifty post-disaster images from the
xView2 dataset~\cite{b22} were streamed directly from the compressed archive
and processed by LLaVA-7B under Ollama without full extraction to disk. A
structured prompt enforced four output fields per image: SUMMARY, DAMAGE,
ACCESS, and FULL\_DESC. If Ollama is unavailable, the pipeline terminates
rather than generating synthetic fallback descriptions, ensuring all imagery
embeddings remain grounded in authentic observations.

\textit{IoT telemetry.} Seven days of 15-minute interval gage-height readings
were downloaded from the USGS NWIS Instantaneous Values API for site~06730200
(Boulder Creek at North 75th Street, Colorado), using parameter code
\texttt{00065}. The 24 most recent measurements were embedded to represent
live conditions; additionally, the top-10 WARNING and CRITICAL readings from
the seven-day window were indexed to preserve recent flood history. Severity
labels (NORMAL, WARNING, DANGER, CRITICAL) were assigned automatically against
the 10.0\,ft flood-stage threshold.

\textit{Official protocols.} Three documents were parsed with PyPDF: the CDC
Preparedness and Response Capabilities manual, the FEMA Disaster Preparedness
Training Manual, and the FEMA Field Operations Guide. Text was segmented using
\texttt{RecursiveCharacterTextSplitter} with $C_{\text{size}}{=}1{,}000$
characters, $C_{\text{overlap}}{=}200$ characters, and a hierarchical
separator sequence (paragraph, newline, sentence, whitespace). The 20\% overlap
prevents procedural steps from being split across chunk boundaries.

\subsection{Live API Integration}

Three unauthenticated public APIs provide live situational intelligence,
improving reproducibility and avoiding dependency on commercial services.
Table~\ref{tab:apis} summarises the endpoints and their operational roles.

\begin{table}[ht]
\centering
\caption{Live API integration summary.}
\label{tab:apis}
\renewcommand{\arraystretch}{1.2}
\begin{tabular}{@{}lll@{}}
\toprule
\textbf{API} & \textbf{Endpoint} & \textbf{Role} \\ \midrule
Open-Meteo   & \texttt{api.open-meteo.com/v1/forecast} & Weather + severity \\
USGS NWIS    & \texttt{waterservices.usgs.gov/nwis/iv/} & River gauge heights \\
USGS Seismic & \texttt{earthquake.usgs.gov/fdsnws/}     & $M{\geq}3.0$ events \\ \bottomrule
\end{tabular}
\end{table}

The weather tool dynamically geocodes location names via the Open-Meteo
geocoding endpoint and derives severity categories (NORMAL to CRITICAL) from
rainfall intensity and wind-speed thresholds. The water-sensor tool computes
the percentage deviation from the site-specific flood-stage threshold and
automatically generates an evacuation warning when readings exceed 130\% of
flood stage. The seismic tool queries a rolling seven-day window filtered to
magnitude $M{\geq}3.0$ within a 300\,km radius of the requested location;
events are ranked by magnitude and annotated with severity labels. All three
tools implement exponential back-off for HTTP~429 rate-limit responses.

\subsection{Software Stack and System Specifications}

Table~\ref{tab:stack} lists the principal software components used throughout
the system. All experiments were conducted on a commodity CPU-based
workstation; GPU acceleration was not required for embedding generation or
text retrieval, though optional GPU acceleration improved LLaVA inference
throughput for the imagery ingestion step.

\begin{table}[ht]
\centering
\caption{Software stack.}
\label{tab:stack}
\renewcommand{\arraystretch}{1.2}
\setlength{\tabcolsep}{4pt}
\begin{tabular}{@{}lll@{}}
\toprule
\textbf{Component}     & \textbf{Library}                    & \textbf{Role} \\ \midrule
Agentic framework      & LangGraph $\geq$0.2                 & ReAct orchestration \\
LLM inference          & LangChain-Groq $\geq$0.2            & Groq API wrapper \\
Embedding model        & \texttt{BAAI/bge-small-en-v1.5}     & 384-dim embeddings \\
Vector store           & ChromaDB $\geq$0.5                  & Persistent index \\
PDF parsing            & PyPDF $\geq$4.0                     & Manual ingestion \\
Vision inference       & Ollama LLaVA-7B                     & Image analysis \\
Operational dashboard  & Streamlit + Plotly                  & Operator UI \\
LLM backbone           & Llama-3.3-70B-Versatile             & Reasoning engine \\ \bottomrule
\end{tabular}
\end{table}

The operational Streamlit dashboard renders tool-call traces as colour-coded
badges (sensor, vision, protocol), maps the agent's STATUS field to the
FEMA/NIMS four-tier colour scheme, and writes every CRITICAL-level response
to a persistent alert log (\texttt{data/alerts.jsonl}) for post-event audit.
Each agent response passes through the safety validator before display,
ensuring no unsafe phrase reaches an operator.
